\section{Actual sections, homogeneous localization and the smooth-projective complex}
\label{sec:actual-section-localization-v054}

\status{This chapter constructs the localization and extension maps for
the actual line-bundle sections on the original scheme.  It proves the
full original twisted-section power-killing and extension theorems,
the dehomogenization and section-cube comparisons, and finite
presentation and projectivity of the actual affine section modules.
Together with FCD18 and GSL8--GSL10, these give the original
\texttt{nonempty\_\allowbreak{}perf\_\allowbreak{}of\_\allowbreak{}smooth\_\allowbreak{}is\allowbreak{}Proj\allowbreak{}By} for every commutative
base ring.  The other polarized, vanishing, rank, multiplication and
square-law suppliers are not conclusions of this chapter.}

The scheme and its bundles remain fixed throughout each construction.
In particular, the section module on an affine chart is not replaced
by an unrelated projective module.  The key geometric argument has two
steps: first clear denominators on finitely many affine charts; then
multiply by a further common power to make those extensions agree on
the overlaps.  Agreement on the nonzero locus alone would not suffice
for gluing over a ring with zero divisors.

\paragraph{Notation and the original coordinate convention.}
For an open $V\subseteq Y$ write $\mathcal O(V)=\Gamma(V,\mathcal O_Y)$.
All restrictions in a displayed identity are understood to the indicated
common open.  A line bundle $M$ is described by opens $U_p$ and units
$a_{pq}$, with coordinate convention $z_p=a_{pq}z_q$.
Write $V_l$ and $b_{lk}$ for the corresponding data of a second
line bundle $H$.  For every $t\in\mathbb Z$, put
\[
 M(t)=M\otimes H^{\otimes t},\qquad
 z_{pl}=a_{pq}b_{lk}^{t}z_{qk}.                 \tag{ASL.1}
\]
Negative tensor powers mean powers of the dual bundle.  In the source
this is literally \texttt{tw M H t}: the transition of
$\mathrm{LB.div}(M,\mathrm{pw}(H,-t))$ is the unit in (ASL.1).
We write $\operatorname{Sec}(M(t),V)$ for its actual compatible
coordinate families, the source's \texttt{Sec\allowbreak{}On}.  Its $R$-module
structure for $f:Y\to\operatorname{Spec}R$ is restriction through
the actual scalar map
\[
 c_f(V):R\longrightarrow\mathcal O(V).
                                                        \tag{ASL.2}
\]
Thus \texttt{Sec\allowbreak{}On\allowbreak{}R} changes the specified scalar action, not the sections.
Multiplication, restriction, pullback and gluing below are operations on
these families.  Ordinary facts about the structure sheaf of an affine
scheme, finite subcovers, and module localization will be used with
their concrete content spelled out where it matters.  No Noetherian,
reducedness or flat-base-change hypothesis on the original $R$ is imposed.

\subsection{The local calculations and the two common exponents}

% BEGIN ASL1
\paragraph{ASL1: fractions on an affine open and power-killing on a compact open.}
\label{claim:ASL1}
Let $W$ be affine, $u\in\mathcal O(W)$, and let $D_W(u)$ be its
principal open.  Every $z\in\mathcal O(D_W(u))$ has the form
$v/u^n$ with $v\in\mathcal O(W)$ and $n\ge0$.  If $v$ restricts
to zero on $D_W(u)$, some power $u^nv$ is zero on $W$.
More generally the latter conclusion holds when $W$ is any
quasi-compact open of any scheme.

\begin{proof}
For $W=\operatorname{Spec}B$, the definition of the structure sheaf
on the basis of principal opens gives
$\mathcal O(D_W(u))=B[u^{-1}]$, with restriction $v\mapsto v/1$.
Every localization element is a fraction with a power of $u$ as
denominator.  The equality $v/1=0$ means exactly that
$u^nv=0$ for some $n$.  Neither assertion cancels $u$.

For a general quasi-compact $W$, choose finitely many affine opens
$W_1,\ldots,W_r$ covering it.  Restriction of $D_W(u)$ to $W_j$
is $D_{W_j}(u|_{W_j})$.  The affine statement therefore supplies
an exponent $n_j$ killing $v|_{W_j}$.  Choose $N\ge n_j$ for all
$j$.  Then $u^Nv$ is zero on each $W_j$, and hence zero by the
separatedness part of the sheaf axiom.  For the empty open the
function ring is zero and exponent zero works.
\end{proof}
% END ASL1

% BEGIN ASL2
\paragraph{ASL2: nonzero loci, units, and actual twisted coordinates.}
\label{claim:ASL2}
Let $\sigma\in\operatorname{Sec}(H^{\otimes d},Y)$, where
$d\in\mathbb Z$, and write $u_l$ for its coordinates.  Then
$u_l=b_{lk}^{d}u_k$ on overlaps.  Its nonzero locus $D(\sigma)$
satisfies, for every open $W\subseteq V_l$,
\[
 D(\sigma)\cap W=D_W(u_l|_W).                       \tag{ASL.3}
\]
Consequently $u_l$ is a unit on every open contained in both
$V_l$ and $D(\sigma)$.  Multiplication by $\sigma^n$ sends a
section of $M(t)$ to a section of $M(t+nd)$, with coordinates
$u_l^nz_{pl}$.  On $D(\sigma)$ this operation is invertible.
These maps commute with every restriction.

\begin{proof}
Multiplication by the transition unit $b_{lk}^d$ does not change
whether the germ belongs to the maximal ideal of the local ring.
Thus the coordinatewise principal opens agree on overlaps and glue
to exactly (ASL.3).  A regular function whose germ is a unit at every
point has a regular inverse: its local inverses agree wherever both
are defined and glue by the structure sheaf.  This proves the unit
assertion without an integrality assumption.

Using (ASL.1), on a common open one has
\[
 u_l^nz_{pl}
 =a_{pq}b_{lk}^{t+nd}u_k^nz_{qk}.                 \tag{ASL.4}
\]
This is the precise transition law for the new twist.  On
$D(\sigma)$ the inverse uses $u_l^{-n}$ and satisfies the inverse
transition identity.  Restricting functions preserves products,
units and their integer powers, so the identities are compatible
with arbitrary restrictions.  The same reasoning shows that
sections in the refined charts of $M(0)$ are naturally the original
sections of $M$: their coordinates are the restrictions of $z_p$,
and the redundant $H$-index glues away.  This identification is also
restriction-natural.
\end{proof}
% END ASL2

% BEGIN ASL3
\paragraph{ASL3: finite simultaneous trivializing charts and the shift identity.}
\label{claim:ASL3}
If $Y$ is quasi-compact and $M,H$ are line bundles, there is a
finite affine cover $W_\alpha$ with
$W_\alpha\subseteq U_{p_\alpha}\cap V_{l_\alpha}$.
If $Y$ is quasi-separated, all $W_\alpha\cap W_\beta$ are
quasi-compact.  On an overlap write $a=a_{p_\alpha p_\beta}$,
$b=b_{l_\alpha l_\beta}$, and $u_\alpha=b^du_\beta$.
For all $t,d\in\mathbb Z$ and $n\ge0$,
\[
 (ab^{t+nd})u_\beta^n=u_\alpha^n(ab^t).          \tag{ASL.5}
\]

\begin{proof}
At every point choose one chart for each bundle.  Affine opens
form a basis of the scheme topology, so an affine neighborhood
lies in their intersection.  Quasi-compactness extracts a finite
subcover.  Quasi-separatedness says that the intersection of
two quasi-compact opens is quasi-compact; apply it to the affine
members.  Finally $u_\alpha^n=b^{nd}u_\beta^n$, and all factors
commute.  The laws of integer powers of a unit give (ASL.5), also
when $d$ or $t$ is negative.  This identity is what changes the
transition functions when the exponent used to clear denominators
is increased.
\end{proof}
% END ASL3

% BEGIN ASL4
\paragraph{ASL4: a local power kills a section in every original chart.}
\label{claim:ASL4}
Let $a\in\operatorname{Sec}(M(t),Y)$ restrict to zero on
$D(\sigma)$, with $\sigma\in\operatorname{Sec}(H^d,Y)$ and
$t,d\in\mathbb Z$.  For every $x\in Y$ there exist an open
neighborhood $W$ and $n\ge0$ such that, for \emph{every} original
pair of charts $(p,l)$ and every
$O\subseteq W\cap U_p\cap V_l$,
\[
                    (u_l|_O)^n a_{pl}|_O=0.          \tag{ASL.6}
\]
This is the full neighborhood statement
\texttt{exists\_\allowbreak{}nhds\_\allowbreak{}pow\_\allowbreak{}mul\_\allowbreak{}tw\allowbreak{}Coord\_\allowbreak{}eq\_\allowbreak{}zero}.

\begin{proof}
Choose an affine neighborhood $W\subseteq U_q\cap V_k$ of $x$.
The function $a_{qk}|_W$ is zero on
$W\cap D(\sigma)=D_W(u_k|_W)$, by ASL2.  ASL1 gives
$(u_k|_W)^n a_{qk}|_W=0$ for some $n$.
Restrict to the open $O$ in the statement.  The transition law gives there
\[
 (u_l|_O)^n a_{pl}|_O
   =a_{pq}|_O\, b_{lk}|_O^{\,t+nd}
                           (u_k|_O)^n a_{qk}|_O=0.
\]
The exponent is independent of $(p,l)$ and $O$ because it was
chosen on the one neighborhood $W$.  This establishes the full
order of quantifiers, rather than a separate exponent for every
subchart.
\end{proof}
% END ASL4

% BEGIN ASL5
\paragraph{ASL5: one global power kills the actual twisted section.}
\label{claim:ASL5}
If $Y$ is quasi-compact, the hypotheses of ASL4 imply that some
$k\ge0$ satisfies
\[
 \forall p,l\qquad u_l^k a_{pl}=0
                 \quad\hbox{on }U_p\cap V_l.         \tag{ASL.7}
\]
Equivalently the actual section $\sigma^k a$ of $M(t+kd)$ is
zero.  No quasi-separatedness is needed.  This proves the full
original \texttt{exists\_\allowbreak{}pow\_\allowbreak{}mul\_\allowbreak{}tw\allowbreak{}Coord\_\allowbreak{}eq\_\allowbreak{}zero}.

\begin{proof}
Apply ASL4 at every point and choose a finite subcover
$W_1,\ldots,W_r$ of the resulting neighborhoods.  Let $k$ be
the maximum of their exponents.  On
$W_j\cap U_p\cap V_l$, multiply (ASL.6) by the remaining
power of $u_l$.  The result is zero for all $j$.
These opens cover $U_p\cap V_l$, so the sheaf axiom gives
(ASL.7) on the full original overlap.  The same single $k$
works for every $(p,l)$.  If the scheme is empty, take $k=0$.
\end{proof}
% END ASL5

% BEGIN ASL6
\paragraph{ASL6: local extensions with one denominator exponent.}
\label{claim:ASL6}
Let $Y$ be quasi-compact, put $V=D(\sigma)$, and let
$y\in\operatorname{Sec}(M(t),V)$.  On a finite simultaneous
affine cover $W_\alpha$ as in ASL3, there are functions
$z_\alpha\in\mathcal O(W_\alpha)$ and one $N\ge0$ such that
\[
 z_\alpha|_{V\cap W_\alpha}
      =u_\alpha^N y_\alpha.                         \tag{ASL.8}
\]
On $W_\alpha\cap W_\beta$, set
\[
 \tau_{\alpha\beta}=a_{p_\alpha p_\beta}
                  b_{l_\alpha l_\beta}^{\,t+Nd},
 \qquad
 \delta_{\alpha\beta}=z_\alpha-\tau_{\alpha\beta}z_\beta.
                                                        \tag{ASL.9}
\]
Then $\delta_{\alpha\beta}$ vanishes on
$V\cap W_\alpha\cap W_\beta$.

\begin{proof}
On $W_\alpha$, the open $V\cap W_\alpha$ is principal with
function $u_\alpha$.  ASL1 writes the regular function
$y_\alpha$ there as $v_\alpha/u_\alpha^{n_\alpha}$.
Take $N\ge n_\alpha$ for every member of the finite cover and
put $z_\alpha=u_\alpha^{N-n_\alpha}v_\alpha$.
This proves (ASL.8), with its equality on every smaller open
obtained by restriction.

On the part of an overlap contained in $V$, the transition law
for $y$ is $y_\alpha=a_{p_\alpha p_\beta}
b_{l_\alpha l_\beta}^{t}y_\beta$.  Equation (ASL.5) with
$n=N$ therefore gives
$u_\alpha^Ny_\alpha=\tau_{\alpha\beta}
u_\beta^Ny_\beta$.  Substitute (ASL.8) to obtain the required
vanishing.  This argument does not assert that the discrepancy
already vanishes on the whole overlap.
\end{proof}
% END ASL6

% BEGIN ASL7
\paragraph{ASL7: extension after one power, including all transition units.}
\label{claim:ASL7}
Let $Y$ be quasi-compact and quasi-separated, $H,M$ line bundles,
$d,t\in\mathbb Z$, and $\sigma\in\operatorname{Sec}(H^d,Y)$.
For every $y\in\operatorname{Sec}(M(t),D(\sigma))$ there are
$k\ge0$ and an actual global section
$a\in\operatorname{Sec}(M(t+kd),Y)$ such that
\[
                  a|_{D(\sigma)}=\sigma^k y.          \tag{ASL.10}
\]
In each original chart this is exactly
$a_{pl}|_{D(\sigma)\cap U_p\cap V_l}=u_l^ky_{pl}$.
This proves the full original \texttt{exists\_\allowbreak{}tw\_\allowbreak{}extend\_\allowbreak{}pow\_\allowbreak{}mul}.

\begin{proof}
Use ASL6.  Each $W_\alpha\cap W_\beta$ is quasi-compact by
ASL3.  On it, (ASL.3) identifies its intersection with
$D(\sigma)$ as the principal open of $u_\alpha$.
ASL1 applied to the discrepancy gives an exponent
$k_{\alpha\beta}$ with
$u_\alpha^{k_{\alpha\beta}}\delta_{\alpha\beta}=0$.
There are only finitely many ordered pairs, so choose one
$K\ge k_{\alpha\beta}$ for all of them.

Put $w_\alpha=u_\alpha^Kz_\alpha$.  The transition unit for
the final twist is
\[
 \tau'_{\alpha\beta}
   =a_{p_\alpha p_\beta}
            b_{l_\alpha l_\beta}^{\,t+(N+K)d}.
\]
Using (ASL.5), now with the starting twist $t+Nd$, gives on
the entire overlap
\[
 \begin{split}
 w_\alpha-\tau'_{\alpha\beta}w_\beta
  &=u_\alpha^Kz_\alpha
             -u_\alpha^K\tau_{\alpha\beta}z_\beta\\
  &=u_\alpha^K\delta_{\alpha\beta}=0.
 \end{split}                                            \tag{ASL.11}
\]
Thus the local coordinates genuinely glue to a section $a$ of
$M(t+(N+K)d)$ over $Y$.  To describe the gluing without suppressing
the original charts, on $W_\alpha\cap U_p\cap V_l$ define its
$(p,l)$-coordinate to be the transition unit from
$(p_\alpha,l_\alpha)$ to $(p,l)$ times $w_\alpha$.
Equation (ASL.11) and the cocycle identities make these functions
agree for different $\alpha$.  They therefore glue on every
$U_p\cap V_l$, and their resulting coordinates obey (ASL.1).

On $V\cap W_\alpha$, (ASL.8) says
$w_\alpha=u_\alpha^{N+K}y_\alpha$.
Transporting this identity to any $(p,l)$ uses (ASL.5) once
more and gives exactly the coordinate equation in the statement.
The finite cover then proves it on all of
$V\cap U_p\cap V_l$.  Take $k=N+K$.
All exponents can depend on the given section; no uniformity in
$y$ or in the twist has been asserted.  The argument works for
negative $t,d$, nilpotent coordinate functions and the empty scheme.
\end{proof}
% END ASL7

\subsection{The graded module and its actual chart comparisons}

Fix a commutative ring $R$, a morphism
$f:Y\to\operatorname{Spec}R$, and sections
$s_0,\ldots,s_m\in\operatorname{Sec}(H,Y)$.
Set $D_i=D(s_i)$ and, for a finite index set $I$,
\[
 D_I=\bigcap_{i\in I}D_i,\quad D_\varnothing=Y,
 \qquad X_I=\prod_{i\in I}X_i\in R[X_0,\ldots,X_m].
                                                        \tag{ASL.12}
\]
The actual graded module on an open $V$ is
\[
 G(V)=\bigoplus_{t\in\mathbb Z}\operatorname{Sec}(M(t),V),
 \qquad G=G(Y),\qquad S=R[X_0,\ldots,X_m].          \tag{ASL.13}
\]
Its elements have finite support in $t$.  Denote the inclusion of
the $t$th summand by $[z]_t$.  The source's supplied witness
\texttt{hx} only certifies commutation of the multiplication maps
defining this very module; it changes none of the coordinate formulas.

% BEGIN ASL8
\paragraph{ASL8: the product section and the exact polynomial action.}
\label{claim:ASL8}
The operators $X_i[z]_t=[s_i z]_{t+1}$ commute and define an
$S$-module structure on $G(V)$.  All restriction maps are graded
$S$-linear maps.  For every $I$ there is an actual section
\[
 \sigma_I=\prod_{i\in I}s_i\in\operatorname{Sec}(H^{|I|},Y),
 \quad (\sigma_I)_l=\prod_{i\in I}(s_i)_l,
 \quad D(\sigma_I)=D_I.                              \tag{ASL.14}
\]
For every $k\ge0$, multiplication by $X_I^k$ on a homogeneous
summand has exactly the coordinate action of $\sigma_I^k$.
More generally, if $F\in S$ is homogeneous of degree $d$, then
$F[z]_t=[z']_{t+d}$, where
\[
 z'_{pl}=F\bigl(c_f(U_p\cap V_l),
                      (s_0)_l,\ldots,(s_m)_l\bigr)z_{pl}.
                                                        \tag{ASL.15}
\]
All functions in (ASL.15) are restricted to $V\cap U_p\cap V_l$.

\begin{proof}
The coordinate for $s_i(s_jz)$ is
$(s_i)_l(s_j)_lz_{pl}$, which is symmetric in $i,j$.
Coordinatewise equality proves commutation, and the universal property
of the polynomial ring therefore gives the stated action.
Multiplication by a section commutes with restriction by ASL2,
so the direct-sum restriction is graded and $S$-linear.

On overlaps, each factor in $(\sigma_I)_l$ changes by the same
$H$-transition unit.  Their product changes by its $|I|$th power,
which is the transition for $H^{|I|}$.  In a local ring a finite
product is a unit if and only if each factor is a unit.  ASL2
therefore proves the equality of nonzero loci in (ASL.14).
The empty product is the constant section $1$ of $H^0$, with
nonzero locus $Y$.

For a monomial, repeated multiplication gives (ASL.15) directly;
its degree is the total number of factors.  Constants act through
(ASL.2).  Summing the monomial formulas proves the assertion for
every homogeneous polynomial, including zero.  This also proves
the source's \texttt{exists\_\allowbreak{}smul\_\allowbreak{}of\_\allowbreak{}tw\allowbreak{}Coord} assertion with
the prescribed target degree, and specializes to the stated
$X_I^k$ formula.
\end{proof}
% END ASL8

% BEGIN ASL9
\paragraph{ASL9: the original very-ample charts give exactly the needed topology.}
\label{claim:ASL9}
Suppose \texttt{Very\allowbreak{}Ample\allowbreak{}By}$(f,H,(s_i))$ holds.  Thus the
$D_i$ cover $Y$, each $D_i$ is affine, and on $D_i$ there are
the given functions $r_{ij}$ with $s_j=r_{ij}s_i$ and with
$R[X_0,\ldots,X_m]\to\mathcal O(D_i)$, $X_j\mapsto r_{ij}$,
surjective.  Then $Y$ is quasi-compact and quasi-separated,
and every nonempty $D_I$ is affine.

\begin{proof}
The finite cover by affine opens proves quasi-compactness.
On $D_i$, the section $s_i$ is invertible in the line-bundle
sense of ASL2.  Hence
\[
 D_I=D_{D_i}\left(\prod_{j\in I\setminus\{i\}}r_{ij}\right)
                       \qquad(i\in I).                \tag{ASL.16}
\]
This proves affineness of every nonempty intersection, with its
actual inclusion in $D_i$.

For clarity the quasi-separatedness assertion follows already
from this affine cover and its affine pairwise intersections.
Given two quasi-compact opens, refine each by finitely many
principal affine opens lying in the $D_i$.  The intersection of
one member from $D_i$ with one from $D_j$ is a principal open
of the affine $D_i\cap D_j$, obtained by restricting their two
defining functions and inverting their product.  It is therefore
quasi-compact.  The intersection of the original two opens is a
finite union of such intersections.  This is precisely the needed
quasi-separatedness condition.  No separation over $R$, or
Noetherian property of $R$, was inserted.
\end{proof}
% END ASL9

% BEGIN ASL10
\paragraph{ASL10: full twisted-section localization at every finite chart set.}
\label{claim:ASL10}
Under ASL9, for every $I$, every $t\in\mathbb Z$, and every
$a\in\operatorname{Sec}(M(t),Y)$,
\[
 a|_{D_I}=0\quad\Longrightarrow\quad
       \exists k\ge0:\ X_I^k[a]_t=0.                 \tag{ASL.17}
\]
For every $y\in\operatorname{Sec}(M(t),D_I)$ there are
$k\ge0$ and $a\in\operatorname{Sec}(M(t+k|I|),Y)$ such that
\[
             [a|_{D_I}]_{t+k|I|}=X_I^k[y]_t.         \tag{ASL.18}
\]
These prove the whole original
\texttt{twisted\_\allowbreak{}sections\_\allowbreak{}localise}, with the source's given
commuting-operator witnesses and every finite $I$.

\begin{proof}
Apply ASL5 to the actual product section $\sigma_I$ in ASL8,
using quasi-compactness from ASL9.  The resulting coordinate
identity is exactly (ASL.17) by (ASL.15) and injectivity of each
summand into the direct sum.

Apply ASL7 to the same section, using both topological properties
from ASL9.  It gives $a|_{D_I}=\sigma_I^ky$ in precisely the
twist $t+k|I|$.  ASL8 translates this to (ASL.18).  In
particular the original existential target degree $t'$ may be
taken to be this explicit integer.  If $I=\varnothing$, take
$k=0$ and use the identity restriction $Y\to Y$ in both statements.
Thus neither the empty vertex nor negative twists are omitted.
\end{proof}
% END ASL10

% BEGIN ASL11
\paragraph{ASL11: a canonical isomorphism from each homogeneous localization.}
\label{claim:ASL11}
For $I$ finite and $t\in\mathbb Z$, let $(G[X_I^{-1}])_t$
be the degree-$t$ homogeneous localization.  An element can be
represented by
\[
       a/X_I^k,\qquad a\in G_{t+k|I|}.
\]
There is an actual $R$-linear isomorphism
\[
 \gamma_{I,t}:(G[X_I^{-1}])_t
                 \xrightarrow{\ \sim\ }
                       \operatorname{Sec}(M(t),D_I),
 \qquad \gamma_{I,t}(a/X_I^k)=\sigma_I^{-k}(a|_{D_I}).
                                                        \tag{ASL.19}
\]
It is characterized by
$\sigma_I^k\gamma_{I,t}(a/X_I^k)=a|_{D_I}$.

\begin{proof}
Multiplication by every $s_i$ with $i\in I$ is invertible on
the graded section module $G(D_I)$, by ASL2.  Their product
and all its powers are therefore invertible.  The graded
restriction $G\to G(D_I)$ consequently extends uniquely to
localization; its degree-$t$ formula is (ASL.19).

More explicitly, equality of fractions $a/X_I^k=a'/X_I^{k'}$
means that a further power of $X_I$ kills
$X_I^{k'}a-X_I^ka'$.  Restricting to $D_I$, all these powers
are invertible operations on sections, so division gives the same
value in (ASL.19).  This also checks the formula directly in the
presence of zero divisors.  Addition uses a common denominator;
scalar multiplication uses (ASL.2).

If the value of $a/X_I^k$ is zero, then $a|_{D_I}=0$.
ASL10 kills the homogeneous numerator by a power of $X_I$,
so its fraction is zero.  Conversely ASL10 applied to a given
$y$ supplies $a$ and $k$ with $a|_{D_I}=\sigma_I^ky$;
the fraction $a/X_I^k$ maps to $y$.  Thus the displayed map
is injective and surjective.  For $I=\varnothing$ this is the
identity on the actual $t$th global-section summand.  These
are exactly the source's \texttt{hom\allowbreak{}Loc} modules and their
fraction identities; no abstract identification of equal ranks is used.
\end{proof}
% END ASL11

% BEGIN ASL12
\paragraph{ASL12: the full section-cube isomorphism, with its restriction maps.}
\label{claim:ASL12}
For every integer $t$, the maps in (ASL.19) form an isomorphism
of the actual cubes
\[
 \gamma_t:\operatorname{grCube}(G,t)
              \xrightarrow{\ \sim\ }
                  \operatorname{secCube}(f,M(t),(D_i)).
                                                        \tag{ASL.20}
\]
Every vertex, including $\varnothing$, is bijective.  This is
the full original
\texttt{exists\_\allowbreak{}gr\allowbreak{}Cube\_\allowbreak{}twisted\allowbreak{}Sections\_\allowbreak{}hom\_\allowbreak{}sec\allowbreak{}Cube\_\allowbreak{}bijective}.
Its coordinate formula commutes with pullback along every morphism
of the original schemes and bundles, and hence along every
specified base change.

\begin{proof}
Let $I\subseteq J$ and put $c=X_{J\setminus I}$.
The localization restriction sends
$a/X_I^k$ to $c^ka/X_J^k$, since $X_J=cX_I$.
Using the actual product sections, its image under $\gamma_{J,t}$
is
\[
 \sigma_J^{-k}\sigma_{J\setminus I}^{k}a|_{D_J}
       =\bigl(\sigma_I^{-k}a|_{D_I}\bigr)|_{D_J}.
                                                        \tag{ASL.21}
\]
All cancelled factors are units on $D_J$, so this identity
does not require them to be non-zero-divisors on $Y$.
It proves commutation with each original restriction map.
Identity and composition follow either from these formulas or
from the corresponding restriction identities for sections.
The vertexwise bijections of ASL11 therefore give the cube
isomorphism, with the precise maps required by the source.

Pullback of a section restricts its coordinate functions by the
specified scheme morphism.  It preserves products, transition
units and their inverses, and sends $D(s_i)$ to the nonzero
locus of the pulled section.  Applying pullback to (ASL.19)
therefore gives the same formula on the pulled data.  Equivalently,
the natural scalar map from the old graded section module to
the new one makes the localized comparison square commute.
This assertion does \emph{not} say that global sections commute
with an arbitrary nonflat base change.  The affine, nonempty
vertices do commute with it, as proved in ASL17 below; the
empty vertex requires no such assertion for the unaugmented
\v Cech complex.
\end{proof}
% END ASL12

% BEGIN ASL13
\paragraph{ASL13: division by the original section and semilinear dehomogenization.}
\label{claim:ASL13}
Let $V\subseteq D(s_i)$ and suppose functions
$r_j\in\mathcal O(V)$ satisfy the literal coordinate identities
$s_j=r_js_i$.  Then $r_i=1$.  Division by $s_i^t$ gives an
injective $\mathcal O(V)$-linear map
\[
 \operatorname{div}_{i,t}:\operatorname{Sec}(M(t),V)
                         \longrightarrow\operatorname{Sec}(M(0),V),
 \qquad z_{pl}\longmapsto(s_i)_l^{-t}z_{pl}.
                                                        \tag{ASL.22}
\]
In fact it is an isomorphism, with inverse multiplication by
$s_i^t$ on these unit coordinates.  The finite-support sum of
these maps defines a semilinear map
\[
 \varepsilon_V:G(V)\longrightarrow\operatorname{Sec}(M(0),V),
 \qquad \varepsilon_V(Fz)=\pi(F)\varepsilon_V(z),
 \quad \pi=\operatorname{eval}(c_f(V),(r_j)).        \tag{ASL.23}
\]
The degree-zero map is the identity.

\begin{proof}
On a common chart $(s_i)_l=b_{lk}(s_i)_k$, and
$z_{pl}=a_{pq}b_{lk}^{t}z_{qk}$.  Thus
\[
 (s_i)_l^{-t}z_{pl}
       =a_{pq}(s_i)_k^{-t}z_{qk},
\]
the transition law for $M(0)$.  All inverse integer powers
are defined because $V\subseteq D(s_i)$; they commute with
restriction.  The inverse formula, injectivity and the
degree-zero identity follow on each chart and hence globally.
The identity $s_i=r_is_i$ implies $r_i=1$ by cancellation
of the unit coordinate, followed by sheaf equality.

For each homogeneous $z\in G(V)_t$,
\[
 \operatorname{div}_{i,t+1}(s_jz)
               =r_j\operatorname{div}_{i,t}(z).
\]
This proves (ASL.23) for $F=X_j$.  It holds for constants by
the specified scalar map (ASL.2), then for monomials by iteration,
and for polynomials by addition.  Every element of $G(V)$ is a
finite sum of homogeneous elements, so this proves the full
semilinearity statement, not only its homogeneous restriction.
\end{proof}
% END ASL13

% BEGIN ASL14
\paragraph{ASL14: full dehomogenization on the actual very-ample chart.}
\label{claim:ASL14}
Let $B_i=\mathcal O(D_i)$ and
$P_i=\operatorname{Sec}(M(0),D_i)$.  For \emph{every} family
$(r_{ij})_j$ in $B_i$ satisfying $s_j=r_{ij}s_i$ in each original
$H$-chart over $D_i$, the actual map
\[
 \varepsilon_i:G\longrightarrow P_i,
 \qquad \varepsilon_i\left(\sum_t[a_t]_t\right)
                     =\sum_t s_i^{-t}a_t|_{D_i}
                                                        \tag{ASL.24}
\]
is semilinear for the original surjection
$\pi_i:S\to B_i$, $X_j\mapsto r_{ij}$, and has all the
source outputs:
\[
 \begin{gathered}
 r_{ii}=1,\qquad \varepsilon_i\text{ is surjective},\\
 a\in G_t,\ \varepsilon_i(a)=0
       \ \Longrightarrow\ \exists k\ge0:\ X_i^ka=0,\\
 \varepsilon_i([a]_0)=a|_{D_i}.
 \end{gathered}                                         \tag{ASL.25}
\]
Thus \texttt{exists\_\allowbreak{}dehom\_\allowbreak{}chart} holds at its full original
scope.  Moreover the actual induced quotient is
$G/(X_i-1)G\simeq P_i$.

\begin{proof}
The coordinate of $s_i$ is a unit on every $D_i\cap V_l$ by
ASL2.  Hence the ratio identities force $r_{ij}$ there to be
the coordinate of $s_j$ divided by that of $s_i$.  Every such
family therefore agrees on these opens with the family supplied
by \texttt{Very\allowbreak{}Ample\allowbreak{}By}.  These opens cover $D_i$,
so sheaf equality proves agreement globally.  In particular the
evaluation $\pi_i$ is the same surjection for every allowed family;
no choice of that family is an additional hypothesis.

Compose graded restriction $G\to G(D_i)$ with ASL13.
This gives (ASL.24), its semilinearity, $r_{ii}=1$, and the
last identity in (ASL.25).
For $y\in P_i$, apply ASL10 at degree zero and $I=\{i\}$.
It gives $a\in G_k$ with $a|_{D_i}=s_i^ky$; hence
$\varepsilon_i(a)=y$.  This proves surjectivity.
For homogeneous $a\in G_t$ in the kernel, injectivity of
division in (ASL.22) gives $a|_{D_i}=0$.  The first part of
ASL10 supplies the stated power of $X_i$.

Finally these are exactly the hypotheses of the already proved
FCD13 quotient argument: $X_i$ acts as $1$ in $P_i$, the map
is surjective, and its homogeneous kernel is killed by elementwise
powers of $X_i$.  Consequently its full kernel is $(X_i-1)G$.
One can see the last step directly by representing any finite sum
modulo $(X_i-1)G$ in one sufficiently high homogeneous degree;
if its image is zero, a power of $X_i$ kills that representative,
and $X_i$ is $1$ in the quotient.  Thus the quotient isomorphism
uses the same semilinear map, rather than a newly chosen module
isomorphism.
\end{proof}
% END ASL14

\subsection{The actual affine section modules}

% BEGIN ASL15
\paragraph{ASL15: ordinary line-bundle sections localize on every qcqs open.}
\label{claim:ASL15}
Let $M$ be a line bundle on any scheme $Y$, let $V$ be a
quasi-compact quasi-separated open, let $f\in\mathcal O(V)$,
and put $W=D_V(f)$.  Then
\[
 \begin{aligned}
 x|_W=0&\ \Longrightarrow\ \exists n\ge0:\ f^nx=0
                    &&(x\in\operatorname{Sec}(M,V)),\\
 \forall y\in\operatorname{Sec}(M,W)&\quad
 \exists n\ge0,\ x\in\operatorname{Sec}(M,V):\quad
                    x|_W=(f|_W)^ny.
 \end{aligned}                                            \tag{ASL.26}
\]
In particular the actual restriction map induces an isomorphism
\[
 \operatorname{Sec}(M,V)[f^{-1}]
                      \simeq\operatorname{Sec}(M,W),
 \qquad x/f^n\longmapsto(f|_W)^{-n}x|_W.            \tag{ASL.27}
\]
This proves the whole original \texttt{sec\allowbreak{}On\_\allowbreak{}basic\allowbreak{}Open\_\allowbreak{}localization}.

\begin{proof}
Work on the open subscheme $V$.  Take $H$ to be its trivial
line bundle, take $d=t=0$, and regard $f$ as a section of
$H^0$.  The actual sections of $M\otimes H^0$ are identified
with those of $M$ by the restriction-natural coordinate map
of ASL2.  ASL5 gives the first assertion of (ASL.26); ASL7
gives the second, with no change of the bundle because $d=0$.
This use of the trivial bundle changes no function or restriction
on the original open $V$.

The displayed fraction formula is well-defined because $f|_W$
is a unit.  The first assertion of (ASL.26) says that a fraction
in its kernel has numerator killed by a further power of $f$,
and is therefore zero.  The second assertion produces a
preimage for every $y$.  Applying the first assertion to a
difference also proves the exact equality criterion for two
restricted sections.  These are precisely the three localization
conditions used by the source's affine chart map.
\end{proof}
% END ASL15

% BEGIN ASL16
\paragraph{ASL16: finite presentation and projectivity of the actual affine section module.}
\label{claim:ASL16}
For every affine open $V$ of any scheme $Y$ and every actual
line bundle $M$, the module
\[
                    P=\operatorname{Sec}(M,V)
              \quad\hbox{over}\quad B=\mathcal O(V)
\]
is finitely presented and projective.  This proves the full
original \texttt{sec\allowbreak{}On\_\allowbreak{}finite\allowbreak{}Presentation\_\allowbreak{}projective} over
arbitrary affine coordinate rings.

\begin{proof}
Identify $V$ with $\operatorname{Spec}B$ by its canonical affine
identification.  Principal opens form a basis.  Refining the
given bundle charts and using quasi-compactness gives a finite
cover $D(f_i)$ on which $M$ has an actual generator $e_i$.
The dual bundle $M^\vee$ has the inverse transition units;
let $e_i^\vee$ be the dual local generator.  Coordinatewise
evaluation gives $M^\vee\otimes M\to\mathcal O_V$ and identifies
the endomorphism given by $u\otimes v$ with multiplication by
the function $v(u)$.

Apply ASL15 to extend $e_i$ and $e_i^\vee$ after powers of $f_i$.
Thus there are global sections $u_i\in P$ and
$v_i\in\operatorname{Sec}(M^\vee,V)$ and exponents $n_i,m_i$
with
\[
 u_i|_{D(f_i)}=f_i^{n_i}e_i,\qquad
 v_i|_{D(f_i)}=f_i^{m_i}e_i^\vee.
\]
Set $h_i=v_i(u_i)\in B$.  On $D(f_i)$ it equals
$f_i^{n_i+m_i}$.  ASL1 therefore gives $l_i\ge0$ such that
\[
               f_i^{l_i}h_i=f_i^{l_i+n_i+m_i}.       \tag{ASL.28}
\]
Put $\widetilde u_i=f_i^{l_i}u_i$,
$q_i=l_i+n_i+m_i$, and define the actual $B$-linear functional
$\lambda_i:P\to B$ by $\lambda_i(z)=v_i(z)$.
On every bundle chart, evaluation in rank one and (ASL.28) give
\[
                 \widetilde u_i\lambda_i(z)=f_i^{q_i}z
                       \qquad(z\in P).               \tag{ASL.29}
\]
This equality holds globally.  The extra factor $f_i^{l_i}$
was included precisely because equality on $D(f_i)$ alone
need not imply equality on $V$.

The functions $f_i^{q_i}$ generate the unit ideal.  If all
$q_i>0$, a prime containing them would contain every $f_i$,
contrary to the covering property.  If any $q_i=0$, one
generator is already $1$.  Hence choose $b_i\in B$ with
$\sum_i b_if_i^{q_i}=1$.  Summing (ASL.29) yields the
explicit finite dual-basis identity
\[
                    z=\sum_i\widetilde u_i\,
                                      b_i\lambda_i(z).  \tag{ASL.30}
\]
Let $\rho:B^r\to P$ send its standard basis to the
$\widetilde u_i$, and let
$j:P\to B^r$ have coordinates $b_i\lambda_i$.
Equation (ASL.30) is $\rho j=1$.  Thus $P$ is a direct
summand of a finite free module, hence finite projective.
For finite presentation explicitly, $\ker\rho$ is the image
of the idempotent $1-j\rho$ on $B^r$ and is generated by
its finitely many columns.  This is a finite presentation
of $P$ over $B$.

When $V$ is empty, $B$ and $P$ are zero and the conclusion
holds directly.  Otherwise the argument uses neither
Noetherianity of $B$ nor a pre-existing equivalence between
line bundles and projective modules.  It constructs the
splitting for this very section module.
\end{proof}
% END ASL16

% BEGIN ASL17
\paragraph{ASL17: actual affine pullback and the nonempty section vertices.}
\label{claim:ASL17}
For every ring map $B\to B'$ and every line bundle $M$ on
$\operatorname{Spec}B$, the canonical section map is an isomorphism
\[
 B'\otimes_B\operatorname{Sec}(M,\operatorname{Spec}B)
       \xrightarrow{\ \sim\ }
       \operatorname{Sec}(M',\operatorname{Spec}B'),
 \qquad b'\otimes z\longmapsto b'\,\pi^*z.          \tag{ASL.31}
\]
It commutes with the actual restriction maps between affine
opens.  In the very-ample situation, this gives, for every
$R\to R'$ and every nonempty $I$,
\[
 R'\otimes_R\operatorname{Sec}(M(t),D_I)
       \simeq\operatorname{Sec}(M'(t),\pi^{-1}D_I).
                                                        \tag{ASL.32}
\]
These comparisons use the supplied pullback square, not merely
one preferred Cartesian model.

\begin{proof}
Use the finite principal trivializing cover $D(f_i)$ in ASL16.
ASL15 identifies the localization of the actual global-section
module with the free rank-one section module on $D(f_i)$.
The same result on the pulled bundle identifies the target of
(ASL.31), localized at the image of $f_i$, with its pulled
rank-one module.  On these localizations the map sends the
actual generator $e_i$ to its pullback.  It is therefore an
isomorphism.  The images of the $f_i$ still generate the unit
ideal, so their principal opens cover $\operatorname{Spec}B'$.

To justify this last local test algebraically, a module whose
localizations at finitely many unit-ideal generators are zero
is zero: each of its elements is killed by a power of every
generator, and these finitely many powers still generate the
unit ideal.  Apply this to the kernel and cokernel of (ASL.31),
using exactness of localization.  That proves the global
isomorphism without flatness of $B'$ over $B$.
The canonical formula on pure tensors shows directly that
it commutes with restriction and with all local transition units.

For (ASL.32), write $D_I=\operatorname{Spec}B$ using ASL9.
Its pullback is canonically
$\operatorname{Spec}(R'\otimes_RB)$; associativity identifies
the left side with
$(R'\otimes_RB)\otimes_B\operatorname{Sec}(M(t),D_I)$.
Apply (ASL.31).  Any specified pullback square has its unique
comparison isomorphism with this affine model, and the pure
tensor formula transports through that isomorphism.  Hence the
coordinate identity holds for the actual specified square.
In particular these are the section comparisons in GSL9.
No assertion of an isomorphism at the global empty vertex is
needed or made for an arbitrary coefficient change.
\end{proof}
% END ASL17

% BEGIN ASL18
\paragraph{ASL18: the exact scalar map on a smooth or flat affine chart.}
\label{claim:ASL18}
Let $R$ be any commutative ring, $f:Y\to\operatorname{Spec}R$,
and $V\subseteq Y$ affine.  Then
\[
 \begin{split}
 f\text{ smooth}&\ \Longrightarrow\ c_f(V):R\to\mathcal O(V)
                                    \text{ is smooth},\\
 f\text{ flat}&\ \Longrightarrow\ c_f(V):R\to\mathcal O(V)
                                    \text{ is flat}.
 \end{split}                                            \tag{ASL.33}
\]
This is the full original \texttt{cst\_\allowbreak{}smooth\_\allowbreak{}and\_\allowbreak{}flat}.

\begin{proof}
Restrict $f$ to the affine source $V$ and the affine target
$\operatorname{Spec}R$.  The affine-local definition of a
smooth scheme morphism says that the induced homomorphism
\[
 \Gamma(\operatorname{Spec}R,\mathcal O)
                              \longrightarrow\mathcal O(V)
\]
is smooth; the analogous affine-local definition of flatness
says that its target is a flat module over the source.
Equivalently one can restrict the local smooth presentations,
or the stalkwise flat module condition, to these affine charts.
These are the usual scheme-to-ring meanings of the two
properties, with no finite-type assumption added to flatness.

The scalar map actually used in the section module is the
composite of this homomorphism with the canonical isomorphism
$R\simeq\Gamma(\operatorname{Spec}R,\mathcal O)$.
Transporting along a ring isomorphism preserves smoothness
and flatness: the polynomial lifting and finite-presentation
conditions, respectively the exact tensor functor, are simply
identified under restriction of scalars.  This composite is
exactly (ASL.2), including the original target-global-sections
identification.  It proves (ASL.33) for that map.
\end{proof}
% END ASL18

\subsection{The exact join with the completed Noetherian model}

% BEGIN ASL19
\paragraph{ASL19: one Noetherian model for the actual section cube.}
\label{claim:ASL19}
Let $R$ be arbitrary, let $f:Y\to\operatorname{Spec}R$ be
smooth, and fix the original line bundles $H,M$ and the
presentation \texttt{Very\allowbreak{}Ample\allowbreak{}By}$(f,H,(s_i))$.
There exist a Noetherian subring $R_0\subseteq R$, a cube
$C_0$ over $R_0$, and an actual cube map
\[
 \Theta:R\otimes_{R_0}C_0
                       \longrightarrow\operatorname{secCube}(f,M,(D_i))
                                                        \tag{ASL.34}
\]
such that every nonempty vertex of $C_0$ is $R_0$-flat,
all cohomology modules of its unaugmented \v Cech complex
are finite over $R_0$, and $\Theta_I$ is bijective for every
nonempty $I$.  These objects are fixed before any later
coefficient ring is chosen.

\begin{proof}
Form the actual $G$ in (ASL.13).  For chart $i$, use
\[
 B_i=\mathcal O(D_i),\qquad
 P_i=\operatorname{Sec}(M(0),D_i),\qquad
 \pi_i(X_j)=r_{ij},\qquad \varepsilon_i\text{ of (ASL.24)}.
\]
The ring map $R\to B_i$ is the original scalar map and is
smooth by ASL18.  ASL16 proves that $P_i$ is finitely
presented and projective over this same $B_i$.
The surjectivity of $\pi_i$ is part of the given
\texttt{Very\allowbreak{}Ample\allowbreak{}By} record.  ASL14 gives
$\pi_i(X_i)=1$, surjectivity of $\varepsilon_i$, and the
required elementwise $X_i$-power annihilation on its
homogeneous kernel.  Its degree-zero restriction is the
actual restriction of sections.

Thus every hypothesis of the full original algebraic theorem
FCD18 has now been supplied for these specific objects.
Reuse that theorem to obtain one $R_0,C_0$ and
\[
          \theta:R\otimes_{R_0}C_0
                          \longrightarrow\operatorname{grCube}(G,0)
\]
with its three original outputs.  ASL12 gives the actual
cube isomorphism $\gamma_0$ to the sections of $M(0)$.
ASL2 identifies those sections with the sections of $M$,
compatibly with every restriction.  Define $\Theta$ to be
this specific composite of $\theta$, $\gamma_0$ and the
inverse of the $M\to M(0)$ identification.

Flatness and finite cohomology of $C_0$ are the outputs of
FCD18.  For nonempty $I$, every factor of $\Theta_I$ is
bijective.  Their compatibility with restrictions makes
(ASL.34) the actual cube map, rather than a list of unrelated
vertex isomorphisms.  No geometric section comparison or
affine projectivity premise remains unproved in this application
of FCD18.  In particular $R_0\to R$ need not be flat.
\end{proof}
% END ASL19

% BEGIN ASL20
\paragraph{ASL20: the full universally valid finite-projective section complex.}
\label{claim:ASL20}
Under the original hypotheses of ASL19, there are a
nonnegative cochain complex $K$ of $R$-modules and an actual
map
\[
                 \phi:K\longrightarrow\check C((D_i),M)
                                                        \tag{ASL.35}
\]
such that each $K^p$ is finite projective, $K^p=0$ for
$p\ge m+1$, and for \emph{every} ring map $R\to R'$,
$R'\otimes_R\phi$ is a quasi-isomorphism.  The target has
the original alternating restriction differential and is the
source's \texttt{cech\allowbreak{}LB}.  This proves all conclusions of
\texttt{exists\_\allowbreak{}finite\_\allowbreak{}projective\_\allowbreak{}complex\_\allowbreak{}cech\allowbreak{}LB}.

\begin{proof}
Apply the already proved GSL8 to the model and actual cube
map (ASL.34).  Its assumptions are exactly the three
outputs of ASL19.  The unaugmented \v Cech term in degree
$p$ uses the $(p+1)$-element subsets of
$\{0,\ldots,m\}$; hence it uses only nonempty vertices.
GSL8 gives (ASL.35), finite projectivity and the asserted
vanishing bound.  The morphism is the scalar extension of
the one Noetherian complex map, followed by the actual
section comparison already fixed in ASL19.

For any subsequent $R\to R'$, the associativity of scalar
extension identifies this with the same Noetherian model
tensored directly by $R'$ over $R_0$.  The universal
quasi-isomorphism assertion in GSL8 applies, without assuming
either ring map is flat.  ASL17 identifies the tensor-extended
section complex with the actual section complex on every
specified pulled-back scheme: finite direct sums and the
restriction differentials commute with the comparisons.
There is therefore no additional assertion about base change
of global sections hidden in (ASL.35).
\end{proof}
% END ASL20

% BEGIN ASL21
\paragraph{ASL21: the full original smooth-projective \texttt{Perf} producer.}
\label{claim:ASL21}
For every commutative ring $R$, every scheme $Y$ and map
$f:Y\to\operatorname{Spec}R$, if $f$ is smooth and satisfies
the original \texttt{Is\allowbreak{}Proj\allowbreak{}By}, then for every actual line
bundle $M$ the literal record $\mathrm{Perf}(f,M)$ exists.
This is the full original
\texttt{nonempty\_\allowbreak{}perf\_\allowbreak{}of\_\allowbreak{}smooth\_\allowbreak{}is\allowbreak{}Proj\allowbreak{}By}, with arbitrary
later coefficient maps and arbitrary specified pullback squares.

\begin{proof}
Unpack \texttt{Is\allowbreak{}Proj\allowbreak{}By} once to obtain $H,m,(s_i)$ and its
\texttt{Very\allowbreak{}Ample\allowbreak{}By} presentation.  ASL9 supplies the finite
cover with affine nonempty intersections.  ASL20 now supplies
the universally valid finite-projective model for its actual
$M$-section complex.  These are all the inputs to the already
proved GSL10, so its construction applies without a remaining
geometric producer.

To identify precisely the output being reused, GSL10 applies
the fixed finite-matrix construction to $K^0\to K^1\to K^2$.
It chooses finite sizes $a,b,c$ and matrices
\[
 D_0:R^a\to R^b,\qquad D_1:R^b\to R^c,
                         D_1D_0=0.                  \tag{ASL.36}
\]
It refines the cover to $W_{ip}=D_i\cap U_p$, where $U_p$
are the original $M$-charts.  Its fixed functions $u_{k,ip}$
are the coordinates on $W_{ip}$ of the $i$th components
of $\phi^0$ applied to the chosen degree-zero elements.
Thus $a,b,c,D_0,D_1,W,\mathrm{ch},u$ are chosen over the
original $R$ before any coefficient change.

Let $g:R\to R'$ be arbitrary, and let
$\pi:Y'\to Y$, $f':Y'\to\operatorname{Spec}R'$ be any
specified Cartesian square.  The GSL10 matrix comparison,
the universal quasi-isomorphism in ASL20, and the actual
section comparison of ASL17 give a bijection
\[
 E_g:\ker(D_0\otimes_RR')
                       \xrightarrow{\ \sim\ }\operatorname{Sec}(M',Y').
                                                        \tag{ASL.37}
\]
In the original pulled coordinates it satisfies
\[
 E_g(v)|_{\pi^{-1}W_{ip}}
       =\sum_{k=1}^{a}c_{f'}(\pi^{-1}W_{ip})(v_k)
                                          \,\pi^*u_{k,ip}.
                                                        \tag{ASL.38}
\]
Indeed each map used to construct (ASL.37) acts on the
original coordinates and the original restriction maps.
ASL12 and ASL17 verify their naturality, including transport
through the unique isomorphism from the supplied pullback
square to the standard affine tensor-product model.
Consequently (ASL.38) is the printed \texttt{Perf.sec} field,
not just an abstract isomorphism of modules.

If the pulled bundle satisfies the original all-cover
$\mathrm{FirstCohomologyVanishes}(M')$ predicate, GSL9
makes its section \v Cech complex exact in degree one.
ASL20 and the GSL10 finite-matrix comparison then give
\[
 \ker(D_1\otimes_RR')=\operatorname{im}(D_0\otimes_RR').
                                                        \tag{ASL.39}
\]
This is exactly \texttt{Perf.exact}, with its own stated
vanishing antecedent.  No such antecedent was needed to
construct the record.  The refinements cover $Y$ and lie
in their indicated original charts, so the remaining
\texttt{le}, \texttt{cover} and composition fields also hold.
All literal fields of \texttt{Perf} have thus been supplied
in the original quantifier order.
\end{proof}
% END ASL21

\paragraph{Completed geometric supplier and retained boundaries.}
ASL21 removes the smooth-projective perfect-section-complex
producer from the remaining inputs of the GSL section-law route.
The Noetherian-model algebra is reused from FCD18, and its
comparison is now attached to the original sections and bundles.
Neither a field nor a reduced base was substituted for the
original arbitrary ring.  The exponent choices are the actual
finite-cover choices, and the complex and its coordinate functions
are fixed before all later base changes.

The independent GSL inputs concerning polarized base change,
Serre or geometric $H^1$ vanishing, the multiplication map and
its invariants, geometric integrality, principal rank one,
torsion-section decomposition and finite growth, and the square
generators remain.  The original \texttt{Moduli} record and
the retained HUF/PUN foundations are unchanged.  Thus this
chapter supplies a full original \texttt{Perf} producer;
it does not by itself supply global \texttt{Sec\allowbreak{}Base\allowbreak{}Change\allowbreak{}St},
\texttt{Sec\allowbreak{}Base\allowbreak{}Change\allowbreak{}Any\allowbreak{}St}, \texttt{Sec\allowbreak{}Twist\allowbreak{}St}, or a whole
selected RH/GRH obligation.

\subsection{Exact source correspondence and inherited proof inputs}

All primary source references are at commit
\texttt{7cbf7a0c32fca1badf05931879a37c3f8a2b03b4} of the
accepted repository.  In the following list $C$ abbreviates
\texttt{src/\allowbreak{}Development/\allowbreak{}Abelian\allowbreak{}Variety\allowbreak{}Theory/\allowbreak{}Proofs/\allowbreak{}}
\texttt{Coherent\allowbreak{}Cohomology\allowbreak{}Proofs/\allowbreak{}}.  Full files and exact
range hashes are supplied in the component manifest.

\begin{itemize}
\item ASL1--ASL7:
$C$\texttt{Exists\allowbreak{}Nhds\allowbreak{}Pow\allowbreak{}Mul\allowbreak{}Tw\allowbreak{}Coord\allowbreak{}Eq\allowbreak{}Zero.lean},
lines 13--62, and
$C$\texttt{Exists\allowbreak{}Tw\allowbreak{}Extend\allowbreak{}Pow\allowbreak{}Mul.lean}, lines 15--171.
The two denominator exponents, all original chart indices
and the integer twists match these complete source statements.
\item ASL8--ASL10:
$C$\texttt{Exists\allowbreak{}Smul\allowbreak{}Of\allowbreak{}Tw\allowbreak{}Coord.lean}, lines 13--49;
$C$\texttt{Mem\allowbreak{}Nonzero\allowbreak{}Locus\allowbreak{}Iff\allowbreak{}Mem\allowbreak{}Basic\allowbreak{}Open\allowbreak{}Chart.lean},
lines 13--59;
$C$\texttt{Very\allowbreak{}Ample\allowbreak{}Compact\allowbreak{}Quasi\allowbreak{}Separated\allowbreak{}Affine\allowbreak{}Charts.lean},
lines 17--50; and
$C$\texttt{Twisted\allowbreak{}Sections\allowbreak{}Localise\allowbreak{}To\allowbreak{}Charts.lean}, lines 24--50.
\item ASL11--ASL12:
$C$\texttt{Exists\allowbreak{}Hom\allowbreak{}Loc\allowbreak{}Bijective\allowbreak{}Gr\allowbreak{}Sum.lean}, lines 45--162,
and $C$\texttt{Exists\allowbreak{}Twisted\allowbreak{}Sections\allowbreak{}Cube\allowbreak{}Bijective.lean},
lines 21--63.  The proof gives the explicit fraction map
whose characterization is used by those source constructors.
\item ASL13--ASL14:
$C$\texttt{Exists\allowbreak{}Dehomogenisation\allowbreak{}Twisted\allowbreak{}Chart.lean},
lines 22--189.  The final quotient-kernel step reuses FCD13.
\item ASL15--ASL18:
$C$\texttt{Sec\allowbreak{}On\allowbreak{}Basic\allowbreak{}Open\allowbreak{}Localization.lean}, lines 14--21;
$C$\texttt{Sec\allowbreak{}On\allowbreak{}Finite\allowbreak{}Presentation\allowbreak{}Projective.lean},
lines 16--65; and $C$\texttt{Cst\allowbreak{}Smooth\allowbreak{}And\allowbreak{}Flat.lean},
lines 12--20.  ASL16 replaces the source's local
finite-presentation/flatness criterion by a proved finite
dual-basis identity on the same module.
The canonical affine pullback map in ASL17 is the map in
\texttt{Proofs/\allowbreak{}Abelian\allowbreak{}Scheme\allowbreak{}Law\allowbreak{}Proofs/\allowbreak{}}
\texttt{Is\allowbreak{}Localized\allowbreak{}Module\allowbreak{}Res\allowbreak{}Basic\allowbreak{}Open\allowbreak{}Of\allowbreak{}Qcqs.lean},
lines 92--115 and 187--261; its actual restriction and
\v Cech comparisons appear there in lines 294--375.
\item ASL19--ASL20:
$C$\texttt{Exists\allowbreak{}Finite\allowbreak{}Projective\allowbreak{}Cech\allowbreak{}Complex\allowbreak{}Line\allowbreak{}Bundle.lean},
lines 19--43.  These use the already reviewed FCD18 and GSL8;
the actual-section base-change input agrees with GSL9.
\item ASL21:
\texttt{Proofs/\allowbreak{}Abelian\allowbreak{}Scheme\allowbreak{}Law\allowbreak{}Proofs/\allowbreak{}}
\texttt{Nonempty\allowbreak{}Perfect\allowbreak{}Sections\allowbreak{}Complex\allowbreak{}Of\allowbreak{}Smooth\allowbreak{}Proj.lean},
lines 15--22, and the literal \texttt{Perf} contract in
\texttt{Definitions/\allowbreak{}Abelian\allowbreak{}Scheme\allowbreak{}Laws/\allowbreak{}}
\texttt{First\allowbreak{}Cohomology\allowbreak{}Vanishes\allowbreak{}Spec\allowbreak{}Of\allowbreak{}Is\allowbreak{}Local\allowbreak{}Ring.lean},
lines 47--81.  The finite-matrix and record construction is
the unchanged reviewed GSL10.
\end{itemize}
